%% annex/p04/universal-reductions.tex --- proof migrations for v1.1 Part V,
%% ch:universal-reductions.  One 40--60 percent compression pass only;
%% evidence statuses are inherited without upgrade.

\begin{annexitem}{annex:p04-universal.1}{thm:p04-renaming-invariance}{Renaming invariance}
\begin{proof}
The transition equation for \(\mathcal T^\iota\) is the
\(\iota\)-conjugate of the one for \(\mathcal T\), so the rewritten symbol
and moved head correspond positionwise.  Induction on steps gives the whole
computation claim.  Initial tapes correspond because blanks are preserved;
the designated states preserve halting answers, and the preserved digit
symbols preserve valid numerical outputs.
\end{proof}
\end{annexitem}

\begin{annexitem}{annex:p04-universal.2}{thm:p04-enumeration}{Adequacy of the numbering: enumeration theorem}
\begin{proof}
(a) For valid \(e\), the fixed wrapper simulates its decoded machine, returns
the canonical word exactly after valid acceptance, and otherwise loops.  For
invalid \(e\), an immediate-loop machine realizes the everywhere-undefined
function.  Thus every \(\varphi_e\) satisfies
Definition~\ref{def:p04-partial-computable}.

(b) Let \(\mathcal T\) witness partial computability of
\(f:\N^k\rightharpoonup\N\).  Since \(\operatorname{code}_1\) is the identity
and higher codes iterate the pairing bijection, \(\operatorname{code}_k\) is
surjective for \(k\ge1\); the required inputs exhaust the binary numerals.
Since an initial configuration exists only for input-alphabet words, the
definition's demand of an accepting or infinite computation on every such
numeral forces \(0,1\in\mathcal A_{\mathrm{in}}\).

If \(q_0\) is halting, add a fresh state \(q'_0\) with
\(\delta(q'_0,a)=(q_0,a,R)\) for every tape symbol.  This extra step preserves
the tape; final head position and off-prefix cells do not affect output, so
accepted outputs and divergence are unchanged.  The three designated states
are now distinct.  Apply a role-preserving bijection onto the designated state
and symbol slots.  By Theorem~\ref{thm:p04-renaming-invariance}, the normalized
machine has the same behavior on binary inputs.

List its full transition graph lexicographically by numbered key and encode
the resulting tuple \(\vec n\) of length \(r\) as
\(e=\operatorname{pair}_{\N}(r,\operatorname{code}_r(\vec n))\).  This is a
valid code, and its decoded machine accepts with
\(\operatorname{bin}(f(\vec x))\) exactly when \(f(\vec x)\) is defined and
otherwise diverges.  Hence
\(\varphi_e(\vec x)\simeq f(\vec x)\) for every \(\vec x\in\N^k\).
\end{proof}
\end{annexitem}

\begin{annexitem}{annex:p04-universal.3}{thm:p04-universal-computation}{Universal computation}
\sketchnote{The simulator's transition encoding and invariant are omitted.}
Invalid codes diverge.  Otherwise update an encoded configuration from the
decoded finite table, returning a canonical output only after valid acceptance
and looping after rejection or invalid acceptance.  Every finite computation
is reproduced, so \(\mathsf{Univ}(e,x)\simeq\varphi_e(x)\).
\end{annexitem}

\begin{annexitem}{annex:p04-universal.4}{thm:p04-reduction-preservation}{Preservation under many-one reduction}
\begin{proof}
Compose the reduction with a decider or recognizer for \(B\); its
biconditional yields the procedure for \(A\).
\end{proof}
\end{annexitem}

\begin{annexitem}{annex:p04-universal.5}{thm:p04-halting-undecidable}{Undecidability of the self-application domain set}
\begin{proof}
\(\mathsf{Univ}(e,e)\) recognizes \(K\).  If
\(h=\operatorname{char}_K\) were computable, define
\(d(e)=0\) when \(h(e)=0\), undefined otherwise.  The load-bearing occurrence
clause of Theorem~\ref{thm:p04-enumeration} gives \(k\) with
\(\varphi_k\simeq d\).  Each value of \(h(k)\) then contradicts the
corresponding definedness of \(\varphi_k(k)\).
\end{proof}
\end{annexitem}

\begin{annexitem}{annex:p04-universal.6}{thm:p04-k-complement}{The complementary domain set}
\begin{proof}
Both c.e. decide \(K\), contradiction.
\end{proof}
\end{annexitem}

\begin{annexitem}{annex:p04-universal.7}{thm:p04-recursion-theorem}{Kleene recursion theorem}
\sketchnote{The code-level compiler for \(p\) and its invariant are omitted;
universal computation and parameterization are inputs.}
They give total computable \(p\) with
\[
  \varphi_{p(x)}(y)\simeq\varphi_{\varphi_x(x)}(y),
\]
with the right side undefined when \(\varphi_x(x)\) is.  By
Theorem~\ref{thm:p04-enumeration}, choose \(a\) for
\(x\mapsto f(p(x))\); put \(e=p(a)\).  Then
\[
  \varphi_e=\varphi_{p(a)}
  \simeq\varphi_{\varphi_a(a)}
  =\varphi_{f(p(a))}=\varphi_{f(e)}.
\]
\end{annexitem}

\begin{annexitem}{annex:p04-universal.8}{thm:p04-rice}{Rice's theorem}
\begin{proof}
Suppose an everywhere-undefined index is outside the set \(A\),
and choose \(a\in A\).
Theorem~\ref{thm:p04-enumeration} gives \(b\) for
\[
 F(e,x)\simeq
 \begin{cases}
   \varphi_a(x),&\text{if }\varphi_e(e)\text{ returns},\\
   \text{undefined},&\text{otherwise},
 \end{cases}
\]
by simulating \(\mathsf{Univ}(e,e)\) before \(\mathsf{Univ}(a,x)\).
Parameterization makes
\(g(e)=\mathsf{par}_1^1(b,e)\) total computable with
\(\varphi_{g(e)}\simeq F(e,\cdot)\).  If \(e\notin K\), extensionality puts
\(g(e)\notin A\); if \(e\in K\), it puts \(g(e)\in A\).  Thus
\(K\leq_m A\), so \(A\) is undecidable.  If an undefined-function index lies
in \(A\), apply this case to the nontrivial extensional complement.
\end{proof}
\end{annexitem}

\begin{annexitem}{annex:p04-universal.9}{thm:p04-numbering-partial-recursive}{Range of the local numbering}
\begin{proof}
If \(\varphi_e\simeq f\), Theorem~\ref{thm:p04-enumeration}(a) makes \(f\)
partial computable, so the cited Theorem~\ref{thm:p04-tm-partial-recursive}
makes it partial recursive.  Conversely, that cited theorem makes every
partial recursive \(f\) partial computable, and
Theorem~\ref{thm:p04-enumeration}(b) supplies a local index \(e\) with
\(\varphi_e\simeq f\).
\end{proof}
\end{annexitem}

\begin{annexitem}{annex:p04-universal.10}{thm:p04-binary-application-equivalence}{Binary and self-application domain problems}
\begin{proof}
Diagonal pairing gives \(K\leq_mH\).  Theorem~\ref{thm:p04-enumeration}
supplies \(b\) with \(\varphi_b(e,x,y)\simeq\varphi_e(x)\).  For
\(z=\operatorname{pair}_{\N}(e,x)\), put \(q=\mathsf{par}_1^2(b,e,x)\).
Parameterization makes \(z\mapsto q\) total with
\(\varphi_q(q)\simeq\varphi_e(x)\), so \(z\in H\) iff \(q\in K\).  Hence
\(H\leq_mK\), whence \(H\sim_mK\) and \(H\) is undecidable.
\end{proof}
\end{annexitem}
